% GENERATED FILE. Edit canonical lessons, then run npm run generate:books.
% canonical-source-hash: fnv1a64:cc42598ed71cdb2e
% canonical-lessons: TE-C521-balahinata, TE-C521-nirasincu, TE-C521-cikkipovu, TE-C521-lavekku, TE-C521-posakaharam

\chapter{Losing Strength}
\label{ch:te-losing-strength}
\hlchaptermodality{521}

\begin{chapteropening}
\textbf{By the end of this chapter:} \emph{I can say that someone grew weak, thin or heavier, and talk about nutrition.}

Complete the last lesson of chapter 521: I can say that someone grew weak, thin or heavier, and talk about nutrition.
\end{chapteropening}

\section[\texorpdfstring{\te{బలహీనత} (balahīnata)}{బలహీనత (balahīnata)}]{\te{బలహీనత} (balahīnata) — frailty; weakness}

\label{lesson:TE-C521-balahinata}



\begin{quote}
Before the new one: say the Telugu for to blow one's nose, then the Telugu for to pick up.
\end{quote}

\subsection*{You'll want to know: \te{బలహీనత}}
\textbf{\te{బలహీనత}} — \emph{balahīnata} — \textquotedblleft{}frailty; weakness\textquotedblright{}.

\begin{quote}
\te{జ్వరం} \te{తరువాత} \te{కొన్ని} \te{రోజులు} \te{బలహీనత} \te{ఉంటుంది}.
\end{quote}

\begin{quote}
\emph{jvaraṁ taruvāta konni rōjulu balahīnata uṇṭundi.}
\end{quote}

\begin{quote}
After a fever there is weakness for a few days.
\end{quote}

\te{బలహీనంగా}, weak; \te{బలం}, strength, + \te{హీన}, lacking.

\begin{cousinweb}
\textbf{\te{బలహీనత}} (\emph{balahīnata}) is a learned compound built from Sanskrit \emph{bala}, \textquotedblleft{}strength\textquotedblright{} (CDIAL 9161), and \emph{hīna}, \textquotedblleft{}lacking\textquotedblright{} (Monier-Williams), with the abstract ending -\te{త}.
\end{cousinweb}

\begin{grammarlens}[title={What you've built}]
One more word for describing what you lived and felt.
\end{grammarlens}

\subsection*{Your turn}
\begin{itemize}
\raggedright
  \item \emph{Say it:} \emph{balahīnata}
  \item \emph{Say it:} \emph{balahīnata}, once more
  \item \emph{Say it:} \emph{jvaraṁ taruvāta konni rōjulu balahīnata uṇṭundi.} — After a fever there is weakness for a few days.
  \item \emph{Recall:} say the Telugu for to blow one's nose, then the Telugu for to pick up, then say \emph{balahīnata} again
\end{itemize}

\begin{tcolorbox}[breakable,colback=teal!4,colframe=teal!35!black,title={Before you move on}]
What does \te{బలహీనత} mean? (\textquotedblleft{}Frailty; weakness\textquotedblright{}.) Say it once more.
\end{tcolorbox}
\section[\texorpdfstring{\te{నీరసించు} (nīrasiñcu)}{నీరసించు (nīrasiñcu)}]{\te{నీరసించు} (nīrasiñcu) — to grow weak}

\label{lesson:TE-C521-nirasincu}



\begin{quote}
Before the new one: say the Telugu for to pick up, then the Telugu for frailty.
\end{quote}

\subsection*{You'll want to know: \te{నీరసించు}}
\textbf{\te{నీరసించు}} — \emph{nīrasiñcu} — \textquotedblleft{}to grow weak\textquotedblright{}.

\begin{quote}
\te{జ్వరంతో} \te{పాప} \te{బాగా} \te{నీరసించింది}.
\end{quote}

\begin{quote}
\emph{jvarantō pāpa bāgā nīrasiñcindi.}
\end{quote}

\begin{quote}
With the fever the little girl grew very weak.
\end{quote}

Past (I …): \te{నీరసించాను}, \emph{nīrasiñcānu}.

The verb of \te{నీరసం}, weakness.

\begin{cousinweb}
\textbf{\te{నీరసించు}} (\emph{nīrasiñcu}) is built on Sanskrit \emph{nīrasa}, \textquotedblleft{}sapless, tasteless\textquotedblright{} (CDIAL 7556), with the Telugu verb ending -\te{ఇంచు}: to lose one's sap.
\end{cousinweb}

\begin{grammarlens}[title={What you've built}]
One more word for describing what you lived and felt.
\end{grammarlens}

\subsection*{Your turn}
\begin{itemize}
\raggedright
  \item \emph{Say it:} \emph{nīrasiñcu}
  \item \emph{Say it:} \emph{nīrasiñcu}, once more
  \item \emph{Say it:} \emph{jvarantō pāpa bāgā nīrasiñcindi.} — With the fever the little girl grew very weak.
  \item \emph{Recall:} say the Telugu for to pick up, then the Telugu for frailty, then say \emph{nīrasiñcu} again
\end{itemize}

\begin{tcolorbox}[breakable,colback=teal!4,colframe=teal!35!black,title={Before you move on}]
What does \te{నీరసించు} mean? (\textquotedblleft{}To grow weak\textquotedblright{}.) Say it once more.
\end{tcolorbox}
\section[\texorpdfstring{\te{చిక్కిపోవు} (cikkipōvu)}{చిక్కిపోవు (cikkipōvu)}]{\te{చిక్కిపోవు} (cikkipōvu) — to grow thin}

\label{lesson:TE-C521-cikkipovu}



\begin{quote}
Before the new one: say the Telugu for frailty, then the Telugu for to grow weak.
\end{quote}

\subsection*{You'll want to know: \te{చిక్కిపోవు}}
\textbf{\te{చిక్కిపోవు}} — \emph{cikkipōvu} — \textquotedblleft{}to grow thin\textquotedblright{}.

\begin{quote}
\te{జబ్బు} \te{తరువాత} \te{అతను} \te{బాగా} \te{చిక్కిపోయాడు}.
\end{quote}

\begin{quote}
\emph{jabbu taruvāta atanu bāgā cikkipōyāḍu.}
\end{quote}

\begin{quote}
After his illness he grew very thin.
\end{quote}

Past (I …): \te{చిక్కిపోయాను}, \emph{cikkipōyānu}.

\te{చిక్కు}, to grow lean, + \te{పోవు}.

\begin{cousinweb}
\textbf{\te{చిక్కిపోవు}} (\emph{cikkipōvu}) is native Dravidian: DEDR 2503 lists Telugu \emph{cikku}, \textquotedblleft{}to become lean or thin\textquotedblright{}, with Tamil \emph{ciṅku}, \textquotedblleft{}to wane\textquotedblright{}; the verb ends in \emph{pōvu}, \textquotedblleft{}to go\textquotedblright{} (DEDR 4572).
\end{cousinweb}

\begin{grammarlens}[title={What you've built}]
One more word for describing what you lived and felt.
\end{grammarlens}

\subsection*{Your turn}
\begin{itemize}
\raggedright
  \item \emph{Say it:} \emph{cikkipōvu}
  \item \emph{Say it:} \emph{cikkipōvu}, once more
  \item \emph{Say it:} \emph{jabbu taruvāta atanu bāgā cikkipōyāḍu.} — After his illness he grew very thin.
  \item \emph{Recall:} say the Telugu for frailty, then the Telugu for to grow weak, then say \emph{cikkipōvu} again
\end{itemize}

\begin{tcolorbox}[breakable,colback=teal!4,colframe=teal!35!black,title={Before you move on}]
What does \te{చిక్కిపోవు} mean? (\textquotedblleft{}To grow thin\textquotedblright{}.) Say it once more.
\end{tcolorbox}
\section[\texorpdfstring{\te{లావెక్కు} (lāvekku)}{లావెక్కు (lāvekku)}]{\te{లావెక్కు} (lāvekku) — to put on weight}

\label{lesson:TE-C521-lavekku}



\begin{quote}
Before the new one: say the Telugu for to grow weak, then the Telugu for to grow thin.
\end{quote}

\subsection*{You'll want to know: \te{లావెక్కు}}
\textbf{\te{లావెక్కు}} — \emph{lāvekku} — \textquotedblleft{}to put on weight\textquotedblright{}.

\begin{quote}
\te{పెళ్ళి} \te{తరువాత} \te{అన్నయ్య} \te{లావెక్కాడు}.
\end{quote}

\begin{quote}
\emph{peḷḷi taruvāta annayya lāvekkāḍu.}
\end{quote}

\begin{quote}
After his wedding my brother put on weight.
\end{quote}

Past (I …): \te{లావెక్కాను}, \emph{lāvekkānu}.

\te{లావు}, stout, + \te{ఎక్కు}, to rise.

\begin{cousinweb}
\textbf{\te{లావెక్కు}} (\emph{lāvekku}) is a native compound: \emph{lāvu}, \textquotedblleft{}bigness, fatness\textquotedblright{} (DEDR 291), and \emph{ekku}, \textquotedblleft{}to increase, to rise\textquotedblright{} (DEDR 768).
\end{cousinweb}

\begin{grammarlens}[title={What you've built}]
One more word for describing what you lived and felt.
\end{grammarlens}

\subsection*{Your turn}
\begin{itemize}
\raggedright
  \item \emph{Say it:} \emph{lāvekku}
  \item \emph{Say it:} \emph{lāvekku}, once more
  \item \emph{Say it:} \emph{peḷḷi taruvāta annayya lāvekkāḍu.} — After his wedding my brother put on weight.
  \item \emph{Recall:} say the Telugu for to grow weak, then the Telugu for to grow thin, then say \emph{lāvekku} again
\end{itemize}

\begin{tcolorbox}[breakable,colback=teal!4,colframe=teal!35!black,title={Before you move on}]
What does \te{లావెక్కు} mean? (\textquotedblleft{}To put on weight\textquotedblright{}.) Say it once more.
\end{tcolorbox}
\section[\texorpdfstring{\te{పోషకాహారం} (pōṣakāhāraṁ)}{పోషకాహారం (pōṣakāhāraṁ)}]{\te{పోషకాహారం} (pōṣakāhāraṁ) — nutrition; nourishing food}

\label{lesson:TE-C521-posakaharam}



\begin{quote}
Before the new one: say the Telugu for to grow thin, then the Telugu for to put on weight.
\end{quote}

\subsection*{You'll want to know: \te{పోషకాహారం}}
\textbf{\te{పోషకాహారం}} — \emph{pōṣakāhāraṁ} — \textquotedblleft{}nutrition; nourishing food\textquotedblright{}.

\begin{quote}
\te{పిల్లలకు} \te{పోషకాహారం} \te{చాలా} \te{అవసరం}.
\end{quote}

\begin{quote}
\emph{pillalaku pōṣakāhāraṁ cālā avasaraṁ.}
\end{quote}

\begin{quote}
Children need nourishing food.
\end{quote}

\te{పోషక}, nourishing, + \te{ఆహారం}, food.

\begin{cousinweb}
\textbf{\te{పోషకాహారం}} (\emph{pōṣakāhāraṁ}) is a learned compound built from Sanskrit \emph{poṣaka}, \textquotedblleft{}nourishing\textquotedblright{} (Monier-Williams), and \emph{āhāra}, \textquotedblleft{}taking food, food\textquotedblright{} (CDIAL 1544).
\end{cousinweb}

\begin{grammarlens}[title={What you've built}]
One more word for describing what you lived and felt.
\end{grammarlens}

\subsection*{Your turn}
\begin{itemize}
\raggedright
  \item \emph{Say it:} \emph{pōṣakāhāraṁ}
  \item \emph{Say it:} \emph{pōṣakāhāraṁ}, once more
  \item \emph{Say it:} \emph{pillalaku pōṣakāhāraṁ cālā avasaraṁ.} — Children need nourishing food.
  \item \emph{Recall:} say the Telugu for to grow thin, then the Telugu for to put on weight, then say \emph{pōṣakāhāraṁ} again
\end{itemize}

\begin{tcolorbox}[breakable,colback=teal!4,colframe=teal!35!black,title={Before you move on}]
What does \te{పోషకాహారం} mean? (\textquotedblleft{}Nutrition; nourishing food\textquotedblright{}.) Say it once more.
\end{tcolorbox}
